\begin{proofbox}
	\textbf{\textcolor{CProof}{思路提示}}

	证明的核心是先将原递推拆分为奇数项和偶数项的子列递推，再分别对$p\neq1$和$p=1$两种情形采用构造法或直接累加法得到通项。

	\medskip
	\textbf{\textcolor{CProof}{正式推导}}
	\begin{description}[style=nextline, leftmargin=2.8em, labelsep=0.5em, itemsep=0.55em]
		\item[\textbf{步骤一（奇偶子列递推）：}]
			在原递推$a_{n+2}=p a_n+q$中分别令$n=2k-1$和$n=2k$，得
			\[
				a_{2k+1}=p a_{2k-1}+q,
				\qquad
				a_{2k+2}=p a_{2k}+q\quad(k=1,2,\dots).
			\]
			\par\textbf{理由：}
			这说明奇数项子列$\{a_{2k-1}\}$和偶数项子列$\{a_{2k}\}$各自满足一个一阶线性递推，且彼此独立。

		\item[\textbf{步骤二（$p\neq1$时构造常数）：}]
			对于递推$a_{2k+1}=p a_{2k-1}+q$，设常数$t$满足$t=pt+q$，解得$t=\dfrac{q}{1-p}$。
			\par\textbf{理由：}
			若令$b_k=a_{2k-1}-t$，则可消去常数项，化为等比形式$b_{k+1}=p b_k$。

		\item[\textbf{步骤三（等比数列求通项）：}]
			由$b_{k+1}=p b_k$知$\{b_k\}$是公比为$p$的等比数列，首项$b_1=a_1-t$。于是
			\[
				b_k=p^{\,k-1}(a_1-t).
			\]
			将$t=\dfrac{q}{1-p}$回代得
			\[
				a_{2k-1}=p^{\,k-1}a_1+q\frac{p^{\,k-1}-1}{p-1}.
			\]
			同理，令$c_k=a_{2k}-t$，可得$a_{2k}=p^{\,k-1}a_2+q\frac{p^{\,k-1}-1}{p-1}$。
			\par\textbf{理由：}
			通过构造常数，将线性递推化为等比数列，再利用等比数列通解即可得到原子列的通项。

		\item[\textbf{步骤四（$p=1$时等差数列）：}]
			当$p=1$时，递推变为$a_{2k+1}=a_{2k-1}+q$，奇数项子列是公差为$q$的等差数列。同理偶数项子列$a_{2k+2}=a_{2k}+q$也是公差为$q$的等差数列。于是
			\[
				a_{2k-1}=a_1+(k-1)q,
				\qquad
				a_{2k}=a_2+(k-1)q.
			\]
			\par\textbf{理由：}
			直接由递推得到相邻两项之差为常数，符合等差数列定义。

		\item[\textbf{步骤五（综合结论）：}]
			综合步骤三和步骤四，即得定理中的完整结论。
			\par\textbf{理由：}
			两种情形的公式覆盖所有可能情况（$p\neq0$已由递推保证）。
	\end{description}

	\medskip
	\textbf{\textcolor{CProof}{结论回扣}}

	\textbf{定理的奇偶通项公式是处理隔项递推的直接工具，使用时只需判断$p$是否等于$1$，代入相应公式即可。}
\end{proofbox}
